% Part: methods
% Chapter: induction
% Section: structural-induction

\documentclass[../../../include/open-logic-section]{subfiles}

\begin{document}

\olfileid[hi]{mth}{ind}{sti}

\olsection{संरचनात्मक आगमन}

अब तक हमने सभी प्राकृतिक संख्याओं के बारे में परिणाम स्थापित करने के लिए
आगमन का उपयोग किया है। किंतु किसी समान सिद्धांत का उपयोग आगमनात्मक रूप से
परिभाषित समुच्चय के सभी !!{element}ों के बारे में परिणाम सीधे सिद्ध करने
में किया जा सकता है। इसे प्रायः \emph{संरचनात्मक} आगमन कहते हैं, क्योंकि
यह आगमनात्मक रूप से परिभाषित वस्तुओं की संरचना पर निर्भर है।

सामान्यतः आगमनात्मक परिभाषा (a) समुच्चय के ``आरंभिक'' !!{element}ों की
सूची और (b) ऐसी संक्रियाओं की सूची से दी जाती है जो पुराने अवयवों से
समुच्चय के नए !!{element} बनाती हैं। उदाहरणतः सुंदर पदों में आरंभिक वस्तुएँ
अक्षर हैं। हमारे पास केवल एक संक्रिया है:
\begin{align*}
  o(s_1, s_2) = & [s_1 \circ s_2]
\end{align*}
आप स्वयं प्राकृतिक संख्याओं $\Nat$ को भी आगमनात्मक परिभाषा से दिया मान
सकते हैं: आरंभिक वस्तु $0$ है और संक्रिया उत्तरवर्ती फलन $x+1$ है।

आगमनात्मक रूप से परिभाषित समुच्चय के सभी अवयवों के बारे में कुछ सिद्ध करने
के लिए, अर्थात् समुच्चय के प्रत्येक !!{element} में कोई गुण $P$ है, हमें:
\begin{enumerate}
\item सिद्ध करना होगा कि आरंभिक वस्तुओं में $P$ है।
\item प्रत्येक संक्रिया $o$ के लिए सिद्ध करना होगा कि यदि उसके तर्कों में
  $P$ है, तो परिणाम में भी $P$ है।
\end{enumerate}
उदाहरणतः सभी सुंदर पदों के बारे में कुछ सिद्ध करने के लिए हम सिद्ध करेंगे
कि वह सभी अक्षरों के बारे में सत्य है, और यदि वह $s_1$ तथा $s_2$ में अलग-अलग
सत्य हो, तो $[s_1\circ s_2]$ के बारे में भी सत्य है।

\begin{prop}
  किसी भी सुंदर पद $t$ में $[$ की संख्या $]$ की संख्या के बराबर है।
\end{prop}

\begin{proof}
हम संरचनात्मक आगमन का उपयोग करते हैं। सुंदर पद आगमनात्मक रूप से परिभाषित
हैं, जहाँ अक्षर आरंभिक वस्तुएँ हैं और पुराने पदों से नए सुंदर पद बनाने की
संक्रिया $o$ है।
\begin{enumerate}
\item दावा प्रत्येक अक्षर के लिए सत्य है, क्योंकि अकेले अक्षर में $[$ की
  संख्या $0$ और उसमें $]$ की संख्या भी $0$ है।
\item मान लें $s_1$ में $[$ की संख्या $]$ की संख्या के बराबर है और यही
  $s_2$ के लिए सत्य है। $o(s_1,s_2)$, अर्थात् $[s_1\circ s_2]$, में $[$
  की संख्या $s_1$ और $s_2$ में $[$ की संख्याओं तथा एक का योग है।
  $o(s_1,s_2)$ में $]$ की संख्या $s_1$ और $s_2$ में $]$ की संख्याओं तथा
  एक का योग है। अतः $o(s_1,s_2)$ में $[$ की संख्या $]$ की संख्या के बराबर
  है।
\end{enumerate}
\end{proof}

\begin{prob}
  संरचनात्मक आगमन से सिद्ध कीजिए कि कोई सुंदर पद $]$ से आरंभ नहीं होता।
\end{prob}

संरचनात्मक आगमन से एक और प्रमाण दें: चिह्नों की शृंखला $t$ का उचित आरंभिक
खंड कोई ऐसी शृंखला $s$ है जो बाएँ से पढ़ने पर चिह्न-दर-चिह्न $t$ से मेल
खाती है, किंतु $t$ लंबी है। अतः उदाहरणतः $[a\circ{}$, $[a\circ b]$ का उचित
आरंभिक खंड है, किंतु न $[b\circ{}$ (वे दूसरे चिह्न पर भिन्न हैं), न
$[a\circ b]$ (उनकी लंबाई समान है)।

\begin{prop}\ollabel{prop:initial}
  सुंदर पद $t$ के प्रत्येक उचित आरंभिक खंड में $]$ की अपेक्षा अधिक $[$ हैं।
\end{prop}

\begin{proof}
  $t$ पर आगमन से:
  \begin{enumerate}
  \item $t$ अपने-आप में एक अक्षर है: तब $t$ का कोई उचित आरंभिक खंड नहीं।
  \item किन्हीं सुंदर पदों $s_1$ और $s_2$ के लिए $t=[s_1\circ s_2]$। यदि
    $r$, $t$ का उचित आरंभिक खंड है, तो कई संभावनाएँ हैं:
    \begin{enumerate}
    \item $r$ केवल $[$ है: तब $r$ में $]$ की अपेक्षा एक अधिक $[$ है।
    \item $r$, $[r_1$ है, जहाँ $r_1$, $s_1$ का उचित आरंभिक खंड है: चूँकि
      $s_1$ सुंदर पद है, आगमन परिकल्पना से $r_1$ में $]$ की अपेक्षा अधिक
      $[$ हैं और $[r_1$ के लिए भी यही सत्य है।
    \item $r$, $[s_1$ अथवा $[s_1\circ{}$ है: पिछले परिणाम से $s_1$ में
      $[$ और $]$ की संख्याएँ बराबर हैं; अतः $[s_1$ अथवा $[s_1\circ{}$ में
      $[$ की संख्या $]$ की संख्या से एक अधिक है।
    \item $r$, $[s_1\circ r_2$ है, जहाँ $r_2$, $s_2$ का उचित आरंभिक खंड
      है: आगमन परिकल्पना से $r_2$ में $]$ की अपेक्षा अधिक $[$ हैं। पिछले
      परिणाम से $s_1$ में $[$ और $]$ की संख्याएँ बराबर हैं। अतः
      $[s_1\circ r_2$ में $[$ की संख्या $]$ की संख्या से अधिक है।
    \item $r$, $[s_1\circ s_2$ है: पिछले परिणाम से $s_1$ में $[$ और $]$
      की संख्याएँ बराबर हैं, और यही $s_2$ के लिए। अतः $[s_1\circ s_2$ में
      $]$ की अपेक्षा एक अधिक $[$ है।
    \end{enumerate}
  \end{enumerate}
\end{proof}

\end{document}
